% !TEX program = xelatex
% ACL 模板说明的中文阅读版。编译：latexmk -xelatex acl_latex.tex
\documentclass[11pt]{article}

% 将 review 改为 final 可生成最终版本（也称 camera-ready 版本）。
% 改为 preprint 可生成非匿名且带页码的版本。
\usepackage[review]{acl}

% 常用宏包。
\usepackage{times}
\usepackage{latexsym}
% 英文原版使用 T1 字体编码和 UTF-8 输入编码；中文版改用 XeLaTeX。
% 此宏包并非必需，可注释；它能改善版式，通常也能节省空间。
\usepackage{microtype}
% 原版用 inconsolata 改善等宽字体；此处由 chinese.tex 的 OpenType 字体设置替代。
% 在 LaTeX 文档中插入图片需要额外加载宏包。
\usepackage{graphicx}
\usepackage{tikz}
% 中文阅读版的共享排版设置；使用 XeLaTeX 编译。
\usepackage{fontspec}
\usepackage{xeCJK}
\setmainfont{NimbusRoman-Regular.otf}[Path=/usr/share/fonts/opentype/urw-base35/,BoldFont=NimbusRoman-Bold.otf,ItalicFont=NimbusRoman-Italic.otf,BoldItalicFont=NimbusRoman-BoldItalic.otf]
\setsansfont{NimbusSans-Regular.otf}[Path=/usr/share/fonts/opentype/urw-base35/,BoldFont=NimbusSans-Bold.otf,ItalicFont=NimbusSans-Italic.otf,BoldItalicFont=NimbusSans-BoldItalic.otf]
\setmonofont{DejaVu Sans Mono}[Scale=MatchLowercase]
\setCJKmainfont{Noto Serif CJK SC}[BoldFont={Noto Serif CJK SC Bold},ItalicFont={Noto Serif CJK SC}]
\setCJKsansfont{Noto Sans CJK SC}[BoldFont={Noto Sans CJK SC Bold}]
\setCJKmonofont{Noto Sans Mono CJK SC}[BoldFont={Noto Sans Mono CJK SC Bold},Scale=0.9908]
\renewcommand{\abstractname}{摘要}
\renewcommand{\figurename}{图}
\renewcommand{\tablename}{表}
\renewcommand{\appendixname}{附录}
\renewcommand{\refname}{参考文献}
\providecommand{\proofname}{证明}
\renewcommand{\proofname}{证明}
% ACL 样式中的匿名作者栏和参考文献标题采用固定英文文本，在此本地化。
\usepackage{etoolbox}
\makeatletter
\patchcmd{\outauthor}{Anonymous ACL submission}{ACL 匿名投稿}{}{}
\patchcmd{\thebibliography}{References\@mkboth{References}{References}}{参考文献\@mkboth{参考文献}{参考文献}}{}{}
\makeatother
\emergencystretch=1em


% 若标题和作者信息超出预留区域，可取消下一行注释，
% 并将 <dim> 设为不小于 5cm 的长度：
% \setlength\titlebox{<dim>}

\title{*ACL 会议论文集排版说明}

% 作者信息可以采用多种格式。
% 多名作者来自同一机构时：
% \author{作者一 \and ... \and 作者 n \\
%         地址第一行 \\ ... \\ 地址最后一行}
% 若姓名无法在同一行放下，可写为：
%         作者一 \\ {\bf 作者二} \\ ... \\ {\bf 作者 n} \\
% 作者来自不同机构时：
% \author{作者一 \\ 地址第一行 \\ ... \\ 地址最后一行
%         \And ... \And
%         作者 n \\ 地址第一行 \\ ... \\ 地址最后一行}
% 使用 \AND 开始新的一排作者，例如：
% \author{作者一 \\ 地址第一行 \\ ... \\ 地址最后一行
%         \AND
%         作者二 \\ 地址第一行 \\ ... \\ 地址最后一行 \And
%         作者三 \\ 地址第一行 \\ ... \\ 地址最后一行}
\author{第一作者 \\
  单位／地址第一行 \\
  单位／地址第二行 \\
  单位／地址第三行 \\
  \texttt{email@domain} \\\And
  第二作者 \\
  单位／地址第一行 \\
  单位／地址第二行 \\
  单位／地址第三行 \\
  \texttt{email@domain} \\}

% 多作者与单位上标的另一种写法：
%\author{
%  \textbf{第一作者\textsuperscript{1}},
%  \textbf{第二作者\textsuperscript{1,2}},
%  \textbf{第三作者\textsuperscript{1}},
%  \textbf{第四作者\textsuperscript{1}},
%\\
%  \textbf{第五作者\textsuperscript{1,2}},
%  \textbf{第六作者\textsuperscript{1}},
%  \textbf{第七作者\textsuperscript{1}},
%  \textbf{第八作者\textsuperscript{1,2,3,4}},
%\\
%  \textbf{第九作者\textsuperscript{1}},
%  \textbf{第十作者\textsuperscript{1}},
%  \textbf{第十一作者\textsuperscript{1,2,3,4,5}},
%  \textbf{第十二作者\textsuperscript{1}},
%\\
%  \textbf{第十三作者\textsuperscript{3}},
%  \textbf{第十四作者\textsuperscript{2,4}},
%  \textbf{第十五作者\textsuperscript{1}},
%  \textbf{第十六作者\textsuperscript{1}},
%\\
%  \textbf{第十七作者\textsuperscript{4,5}},
%  \textbf{第十八作者\textsuperscript{3,4}},
%  \textbf{第十九作者\textsuperscript{2,5}},
%  \textbf{第二十作者\textsuperscript{1}}
%\\
%\\
%  \textsuperscript{1}单位一,
%  \textsuperscript{2}单位二,
%  \textsuperscript{3}单位三,
%  \textsuperscript{4}单位四,
%  \textsuperscript{5}单位五
%\\
%  \small{\textbf{通信作者：} \href{mailto:email@domain}{email@domain}}
%}

\begin{document}
\maketitle
\begin{abstract}
本文档是 *ACL 作者通用指南的补充，说明如何使用 ACL 会议的 \LaTeX{} 样式文件。
文档本身遵循其所述规范，因此也展示了稿件应有的排版效果。
这些说明同时适用于送审论文和录用论文的最终版本。
\end{abstract}

\section{引言}

这些说明面向使用 \LaTeX{} 向 *ACL 会议投稿的作者，并不能替代全部投稿要求。所有作者均须遵守 *ACL 论文集的通用指南，\footnote{\url{http://acl-org.github.io/ACLPUB/formatting.html}} 本文仅补充说明 \LaTeX{} 样式文件的使用方式。

模板包括本文档的 \LaTeX{} 源文件（\texttt{acl\_latex.tex}）、
用于排版的 \LaTeX{} 样式文件（\texttt{acl.sty}）、
ACL 参考文献样式（\texttt{acl\_natbib.bst}）、
示例参考文献库（\texttt{custom.bib}），
以及 ACL Anthology 的参考文献库（\texttt{anthology.bib}）。

\section{编译引擎}

生成 PDF 文件时，强烈建议使用 pdf\LaTeX{}，而非原始 \LaTeX{} 加 dvips+ps2pdf 或 dvipdf。
样式文件 \texttt{acl.sty} 也支持 lua\LaTeX{} 和 Xe\LaTeX{}；这两者尤其适合非拉丁文字。
仓库中的 \texttt{acl\_lualatex.tex} 展示了如何在 lua\LaTeX{} 或 Xe\LaTeX{} 下使用 \texttt{acl.sty}。
本中文阅读版使用 Xe\LaTeX{}。

\section{导言区}

文件第一行必须是
\begin{quote}
\begin{verbatim}
\documentclass[11pt]{article}
\end{verbatim}
\end{quote}

送审版本通过以下命令加载样式文件：
\begin{quote}
\begin{verbatim}
\usepackage[review]{acl}
\end{verbatim}
\end{quote}
最终版本则省略 \verb|review| 选项：
\begin{quote}
\begin{verbatim}
\usepackage{acl}
\end{verbatim}
\end{quote}

若要使用 Times Roman，请在导言区加入：
\begin{quote}
\begin{verbatim}
\usepackage{times}
\end{verbatim}
\end{quote}
也可以使用 txfonts 或 newtx 等替代方案。

其他可能有用的宏包说明，参见本文档 \LaTeX{} 源文件中的注释。

使用 \verb|\title| 和 \verb|\author| 设置标题与作者。作者列表中可用 \verb|\and|、\verb|\And| 和 \verb|\AND| 排列多名作者；示例见 \LaTeX{} 源文件。

默认情况下，标题和作者姓名所在区域的高度设为最小值 5 cm。若需要更多空间，请在导言区加入：
\begin{quote}
\begin{verbatim}
\setlength\titlebox{<dim>}
\end{verbatim}
\end{quote}
将 \verb|<dim>| 替换为所需长度；该长度不得小于 5 cm。

\section{文档正文}

\subsection{脚注}

使用 \verb|\footnote| 命令插入脚注。\footnote{这是一个脚注。}

\subsection{表格与图片}

表~\ref{tab:accents} 展示了表格及表题的示例。
\textbf{不要修改默认的图表标题字号。}

\begin{table}
  \centering
  \begin{tabular}{lc}
    \hline
    \textbf{命令} & \textbf{输出} \\
    \hline
    \verb|{\"a}|     & {\"a}           \\
    \verb|{\^e}|     & {\^e}           \\
    \verb|{\`i}|     & {\`i}           \\
    \verb|{\.I}|     & {\.I}           \\
    \verb|{\o}|      & {\o}            \\
    \verb|{\'u}|     & {\'u}           \\
    \verb|{\aa}|     & {\aa}           \\\hline
  \end{tabular}
  \begin{tabular}{lc}
    \hline
    \textbf{命令} & \textbf{输出} \\
    \hline
    \verb|{\c c}|    & {\c c}          \\
    \verb|{\u g}|    & {\u g}          \\
    \verb|{\l}|      & {\l}            \\
    \verb|{\~n}|     & {\~n}           \\
    \verb|{\H o}|    & {\H o}          \\
    \verb|{\v r}|    & {\v r}          \\
    \verb|{\ss}|     & {\ss}           \\
    \hline
  \end{tabular}
  \caption{带重音符号字符的命令示例，可用于 Bib\TeX{} 条目等场景。}
  \label{tab:accents}
\end{table}

图中的字体应尽量与正文字体一致。图~\ref{fig:experiments} 展示了图片及图注的示例。

使用 \verb|graphicx| 宏包，可以在正文适当位置的 figure 环境中插入图像文件。
该宏包提供多种可选参数，用于控制图片外观。
必须在 \LaTeX{} 导言区中，即 \verb|\documentclass| 声明之后、\verb|\begin{document}| 之前，通过 \verb|\usepackage{graphicx}| 显式加载它。

\begin{figure}[t]
  % 以矢量示例图替代英文占位图，保留黄金矩形的比例。
  \begin{tikzpicture}[x=\columnwidth,y=\columnwidth]
    \fill[gray!30] (0,0) rectangle (1,0.618);
    \draw[gray] (0,0) -- (1,0.618) (0,0.618) -- (1,0);
    \draw (0,0) rectangle (1,0.618);
    \node[font=\Large] at (0.5,0.32) {黄金比例};
    \node[font=\scriptsize] at (0.5,0.09) {原始尺寸：$32.36\times20$ bp};
  \end{tikzpicture}
  \caption{这是一条超过一行的图注。
    示例图片通常由 \texttt{mwe} 宏包提供，即使不在导言区显式加载该宏包也可使用。}
  \label{fig:experiments}
\end{figure}

\begin{figure*}[t]
  \includegraphics[width=0.48\linewidth]{example-image-a} \hfill
  \includegraphics[width=0.48\linewidth]{example-image-b}
  \caption{展示如何并排放置两张图片的最小可运行示例。}
\end{figure*}

\subsection{超链接}

较旧版本的 \LaTeX{} 在编译时可能遇到以下错误：
\begin{quote}
\verb|\pdfendlink| 与 \verb|\pdfstartlink| 的嵌套层级不同。
\end{quote}
使用 pdf\LaTeX{} 且一条引文跨越分页边界时，会出现该问题。最佳解决办法是将 \LaTeX{} 升级到 2018-12-01 或更高版本。

\subsection{引文}

\begin{table*}
  \centering
  \begin{tabular}{lll}
    \hline
    \textbf{输出}           & \textbf{natbib 命令} & \textbf{ACL 专用命令} \\
    \hline
    \citep{Gusfield:97}       & \verb|\citep|           &                           \\
    \citealp{Gusfield:97}     & \verb|\citealp|         &                           \\
    \citet{Gusfield:97}       & \verb|\citet|           &                           \\
    \citeyearpar{Gusfield:97} & \verb|\citeyearpar|     &                           \\
    \citeposs{Gusfield:97}    &                         & \verb|\citeposs|          \\
    \hline
  \end{tabular}
  \caption{\label{citation-guide}
    样式文件支持的引文命令。
    该样式基于 natbib 宏包，支持其全部引文命令。
    为保持兼容，也支持早期 ACL 样式文件中定义的命令。
  }
\end{table*}

表~\ref{citation-guide} 展示了样式文件支持的语法。
我们建议使用 natbib 风格。
\verb|\citet|（正文中的引文）生成“作者（年份）”形式，例如对 \citet{Gusfield:97} 的引用。
\verb|\citep|（括号内引文）生成“（作者，年份）”形式~\citep{Gusfield:97}。
\verb|\citealp|（不带括号的替代引文）生成“作者，年份”形式，适合在已有括号中引用文献（例如 \citealp{Gusfield:97}）。

\verb|\citeposs| 可以生成所有格形式的引文。
它不是标准 natbib 命令，因此通常与其他样式文件不兼容。

\subsection{参考文献}

\nocite{Ando2005,andrew2007scalable,rasooli-tetrault-2015}

提供的 \LaTeX{} 和 Bib\TeX{} 样式文件大体遵循美国心理学会（APA）的格式。
若自建文献库命名为 \texttt{custom.bib}，则在 \LaTeX{} 文件的所有附录之前加入以下内容，即可生成参考文献部分：
\begin{quote}
\begin{verbatim}
\bibliography{custom}
\end{verbatim}
\end{quote}

完整 ACL Anthology 的 Bib\TeX{} 文献库可从 \url{https://aclweb.org/anthology/anthology.bib.gz} 获取。
若要同时使用 Anthology 和自建 .bib 文件，请将上述命令改为：
\begin{quote}
\begin{verbatim}
\bibliography{anthology,custom}
\end{verbatim}
\end{quote}

关于如何准备 Bib\TeX{} 文件，请参见第~\ref{sec:bibtex} 节。

\subsection{公式}

下面给出一个公式示例：
\begin{equation}
  \label{eq:example}
  A = \pi r^2
\end{equation}

公式、节、小节、图和表的编号标签均用 \verb|\label{label}| 定义，再用 \verb|\ref{label}| 进行交叉引用。

这里是对式~\ref{eq:example} 的交叉引用示例。

\subsection{附录}

在任何附录章节之前使用 \verb|\appendix|，可将章节编号切换为字母。示例见附录~\ref{sec:appendix}。

\section{Bib\TeX{} 文件}
\label{sec:bibtex}

原版说明要求不要在 Bib\TeX{} 条目中使用 Unicode；某些特殊字符的输入方式可能干扰 Bib\TeX{} 的字母排序。推荐写法见表~\ref{tab:accents}。
本中文阅读版在 Xe\LaTeX{} 下显示文献标题的中文译名，保留拉丁字母作者信息以便排序，并将英文原题保存在文献库的 \texttt{origtitle} 字段中。

请尽可能在 Bib\TeX{} 记录中提供 DOI 或 URL，并为引用的所有 ACL 文献提供这些信息。
DOI 使用 \verb|doi| 字段，URL 使用 \verb|url| 字段。
若条目包含 URL 或 DOI 字段，参考文献中的论文标题会借助 hyperref \LaTeX{} 宏包成为指向论文的超链接。

\section*{局限性}

本文档不涵盖 ACL 或其他具体会议对内容的要求。
最大页数、必须包含的“局限性”部分等信息，请查阅相应作者指南。

\section*{致谢}

本文档由 Steven Bethard、Ryan Cotterell 和 Rui Yan 在早期 ACL 与 NAACL 论文集说明的基础上改编而成，所参考的说明包括：
Douwe Kiela 和 Ivan Vuli\'{c} 编写的 ACL 2019 说明，
Stephanie Lukin 和 Alla Roskovskaya 编写的 NAACL 2019 说明，
Shay Cohen、Kevin Gimpel 和 Wei Lu 编写的 ACL 2018 说明，
Margaret Mitchell 和 Stephanie Lukin 编写的 NAACL 2018 说明，
Jason Eisner 提供的 (NA)ACL 2017/2018 Bib\TeX{} 建议，
Dan Gildea 和 Min-Yen Kan 编写的 ACL 2017 说明，
Margaret Mitchell 编写的 NAACL 2017 说明，
Maggie Li 和 Michael White 编写的 ACL 2012 说明，
Jing-Shin Chang 和 Philipp Koehn 编写的 ACL 2010 说明，
Johanna D. Moore、Simone Teufel、James Allan 和 Sadaoki Furui 编写的 ACL 2008 说明，
Hwee Tou Ng 和 Kemal Oflazer 编写的 ACL 2005 说明，
Eugene Charniak 和 Dekang Lin 编写的 ACL 2002 说明，
以及 John Chen、Henry S. Thompson、Donald Walker 等人编写的更早 ACL 与 EACL 格式。
另有部分内容取自国际人工智能联合会议（IJCAI）和计算机视觉与模式识别会议（CVPR）的排版指南。

% 同时加载 Anthology 与自建参考文献库的示例：
%\bibliography{custom,anthology-overleaf-1,anthology-overleaf-2}
% 此处仅加载自建参考文献库。
\bibliography{custom}

\appendix
\section{附录示例}
\label{sec:appendix}

这是一个附录。

\end{document}
